% ============================================================
%  CHUONG 4. KET QUA NGHIEN CUU
% ============================================================
\chapter{KẾT QUẢ NGHIÊN CỨU}
\label{ch:ket-qua}

\section{Môi trường triển khai và dữ liệu thử nghiệm}
\label{sec:moi-truong}

Hệ thống được chạy thử trên nền tảng .NET 10 với ASP.NET Core MVC, cơ sở dữ liệu
SQL Server 2022 chạy trong container, trình duyệt kiểm thử là Google Chrome ở
hai độ phân giải 1440 $\times$ 900 và 390 $\times$ 844. Toàn bộ ảnh chụp màn
hình trong chương này được lấy trực tiếp từ hệ thống đang chạy trên môi trường
đó, không qua chỉnh sửa nội dung.

Dữ liệu thử nghiệm do \code{DbInitializer} nạp tự động ở lần chạy đầu tiên trên
cơ sở dữ liệu trống, gồm \textbf{26 món ăn}, \textbf{40 nguyên liệu},
\textbf{10 danh mục} và một tài khoản quản trị mặc định.

\section{Giao diện các chức năng}
\label{sec:giao-dien}

\subsection{Trang chủ và duyệt món ăn}

Trang chủ giới thiệu các danh mục, món nổi bật và số liệu tổng quan của kho dữ
liệu (chức năng A1), từ đó dẫn sang trang danh sách món ăn.

\begin{figure}[H]
\centering
\includegraphics[width=0.65\textwidth]{h4-02-danh-sach-mon.jpg}
\caption{Trang danh sách món ăn với bộ lọc bên trái và lưới kết quả}
\label{fig:h4-02}
\end{figure}

Màn hình danh sách gồm ô tìm kiếm theo tên (chức năng A3), nhóm bộ lọc bên trái
theo danh mục, vùng miền, độ khó và thời gian nấu tối đa (A4), hộp chọn sắp xếp
với bốn phương án tên A-Z, nấu nhanh, calo thấp, calo cao (A5), và lưới thẻ món
phân trang 9 món mỗi trang; các liên kết phân trang giữ nguyên trạng thái lọc và
sắp xếp hiện tại. Nút thêm yêu thích trên từng thẻ chỉ hiện với người dùng đã
đăng nhập. Bấm vào một thẻ mở trang chi tiết, nơi hiển thị đầy đủ thông số, danh
sách nguyên liệu kèm định lượng và các bước nấu.

\subsection{Tìm kiếm, lọc và sắp xếp}

Khi lọc theo vùng miền Nam, độ khó Dễ và sắp xếp theo năng lượng tăng dần, số
món giảm từ 26 xuống còn tập con thỏa mãn cả hai điều kiện, đồng thời thứ tự
hiển thị tuân theo tiêu chí đã chọn; các liên kết sắp xếp và phân trang đều giữ
nguyên trạng thái lọc hiện tại. Tính đúng đắn của từng phương án sắp xếp được
kiểm chứng bằng các ca kiểm thử TC-24 tới TC-29 ở mục~\ref{sec:kiem-thu}.

\subsection{Gợi ý theo nguyên liệu}

Trang gợi ý liệt kê toàn bộ nguyên liệu dưới dạng bảng hộp kiểm sắp theo tên.
Khi người dùng chưa chọn nguyên liệu nào, trang giữ trạng thái ban đầu kèm hướng
dẫn sử dụng thay vì trả về danh sách rỗng.

\begin{figure}[H]
\centering
\includegraphics[width=0.67\textwidth]{h4-06-goi-y-ket-qua.jpg}
\caption{Kết quả gợi ý với 8 nguyên liệu đầu vào}
\label{fig:h4-06}
\end{figure}

Với đầu vào gồm 8 nguyên liệu (chanh, dầu ăn, đường, hành lá, muối, nước mắm, rau
muống, tỏi), hệ thống trả về \textbf{24 món phù hợp}, trong đó \textbf{2 món nấu
được ngay} và 22 món còn thiếu nguyên liệu. Hai món nấu được ngay được xếp lên
đầu danh sách và mang nhãn màu xanh ``Đủ nguyên liệu''; các món còn lại mang nhãn
màu hổ phách ghi rõ số nguyên liệu còn thiếu, kèm danh sách tên nguyên liệu đó
ngay dưới tên món. Nhãn ``Đủ nguyên liệu'' được gắn cho món có độ phủ bằng 1,
tương ứng $M = \emptyset$; thứ tự các thẻ kết quả phản ánh trực tiếp bốn tiêu
chí xếp hạng đã thiết kế ở mục~\ref{sec:thuat-toan-goi-y}.

\subsection{Thực đơn tuần}

Trang thực đơn gồm biểu mẫu sinh thực đơn (nhập tên, chọn tuần bắt đầu và vùng
miền nếu muốn, ứng với chức năng B1) và danh sách các thực đơn đã lưu của riêng
người dùng đang đăng nhập (B4).

\begin{figure}[H]
\centering
\includegraphics[width=0.70\textwidth]{h4-08-thuc-don-tuan.jpg}
\caption{Lịch thực đơn 7 ngày $\times$ 3 bữa}
\label{fig:h4-08}
\end{figure}

Lưới thực đơn hiển thị đủ 21 suất ăn, mỗi ô ứng với một bản ghi
\code{MenuPlanItem} gồm tên món, năng lượng theo khẩu phần, một hộp chọn để đổi
sang món khác và một nút xóa (chức năng B3); cột của ngày vẫn được giữ lại kể cả
khi cả ba bữa đều trống. Dòng cuối cùng hiển thị tổng năng lượng theo từng ngày,
lần lượt là 1580, 1270, 1130, 1270, 1110, 1030 và 820 kcal, được tính lại sau
mỗi lần chỉnh sửa. Nút ``Tạo danh sách đi chợ'' chuyển sang chức năng B5 với mã
thực đơn hiện tại.

\subsection{Danh sách đi chợ}

\begin{figure}[H]
\centering
\includegraphics[height=0.44\textheight]{h4-09-di-cho.jpg}
\caption{Danh sách đi chợ sinh từ thực đơn ở Hình~\ref{fig:h4-08}}
\label{fig:h4-09}
\end{figure}

Từ thực đơn 21 suất ăn, hệ thống sinh ra danh sách gồm \textbf{35 dòng nguyên
liệu} được gộp theo \textbf{8 nhóm}: Đậu/Đỗ, Gia vị, Hải sản, Nấm, Ngũ cốc/Tinh
bột, Rau củ, Thịt và Trứng/Sữa; tiêu đề nhóm lấy theo cột \code{Group} của bảng
\code{Ingredients}. Mỗi dòng hiển thị tên nguyên liệu, tổng khối lượng đã cộng
dồn và đơn vị, kèm hộp kiểm ``đã mua'' đảo trạng thái \code{IsPurchased}; thanh
tiến độ ở đầu trang cập nhật theo số dòng đã đánh dấu. Nút tạo lại danh sách ghi
đè danh sách cũ để phản ánh thực đơn vừa chỉnh sửa.

\subsection{Tài khoản và phân quyền}

Biểu mẫu đăng nhập xác thực bằng thư điện tử và mật khẩu; khi thông tin sai, hệ
thống báo lỗi chung mà không tiết lộ sai ở trường nào.

Cơ chế kiểm tra hợp lệ của biểu mẫu đăng ký hoạt động đồng thời trên ba trường:
sai định dạng email, mật khẩu ngắn hơn 8 ký tự và mật khẩu xác nhận không khớp.
Thông báo lỗi hiển thị bằng tiếng Việt ngay dưới ô nhập tương ứng, viền ô nhập
chuyển sang màu đỏ.

Yêu cầu phi chức năng N2 cũng được kiểm chứng trực tiếp trên hệ thống đang chạy:
một tài khoản vừa đăng ký với vai trò \code{User} khi truy cập đường dẫn
\code{/Admin/Recipes} bị chuyển hướng sang trang thông báo không đủ quyền, kèm
tham số \code{ReturnUrl} ghi nhận trang đã bị chặn.

Sau khi đăng nhập, người dùng còn sử dụng được trang danh sách món yêu thích
(chức năng B7): thêm hoặc bỏ yêu thích ngay trên thẻ món và xem lại toàn bộ món
đã đánh dấu; các món này được thuật toán sinh thực đơn ưu tiên chọn.

\subsection{Khu vực quản trị}

Khu vực quản trị gồm bảng điều khiển với thẻ số liệu tổng quan (chức năng C1),
các bảng quản lý món ăn, nguyên liệu và danh mục với nút thêm, sửa, xóa (C2 tới
C4). Biểu mẫu món ăn cho phép nhập thông số món rồi gán từng nguyên liệu kèm
khối lượng và đơn vị ngay trên cùng một trang. Thao tác xóa nguyên liệu đang
được tham chiếu hoặc danh mục còn món sử dụng bị chặn kèm thông báo giải thích,
đúng như thiết kế hành vi xóa ở mục~\ref{sec:hanh-vi-xoa}.

\subsection{Giao diện đáp ứng và chế độ tối}

Trên màn hình hẹp, bố cục các trang tự sắp xếp lại thành một cột và nhóm bộ lọc
của trang danh sách chuyển thành ngăn kéo trượt mở từ cạnh màn hình. Hệ thống
cũng hỗ trợ chế độ tối: toàn bộ cặp màu chữ và
nền được định nghĩa lại qua biến CSS, và các trang hiển thị trạng thái rỗng có
hướng dẫn thay vì bảng trống khi chưa có dữ liệu.

\section{Kịch bản demo}
\label{sec:demo}

Luồng sử dụng xuyên suốt ba chức năng lõi diễn ra như sau. Người dùng mở mục
``Gợi ý nguyên liệu'', tick 8 nguyên liệu đang có rồi bấm ``Gợi ý món'' và nhận
về 24 món xếp hạng với 2 món nấu được ngay ở đầu danh sách
(Hình~\ref{fig:h4-06}); mở chi tiết vài món và bấm ``Thêm yêu thích'' để các món
đó được ưu tiên khi sinh thực đơn. Sang mục ``Thực đơn tuần'', nhập tên, chọn
tuần bắt đầu rồi bấm ``Tạo thực đơn'', hệ thống sinh đủ 21 suất ăn trong một
thao tác (Hình~\ref{fig:h4-08}); ô nào chưa vừa ý thì chọn món khác ngay trong ô
đó và tổng năng lượng của ngày được tính lại. Cuối cùng bấm ``Tạo danh sách'' để
gộp nguyên liệu của 21 suất thành 35 dòng theo 8 nhóm (Hình~\ref{fig:h4-09});
khi đi chợ, mua xong nguyên liệu nào thì tick vào ô tương ứng, dòng đó bị gạch
ngang và thanh tiến độ tăng lên.

\section{Kết quả kiểm thử}
\label{sec:kiem-thu}

\subsection{Phương pháp kiểm thử}

Kiểm thử được thực hiện ở mức \textbf{kiểm thử đơn vị} bằng thư viện xUnit, tập
trung vào ba lớp dịch vụ chứa thuật toán lõi cùng chức năng sắp xếp danh sách
món. Cơ sở dữ liệu được thay bằng nhà cung cấp \textbf{EF Core In-Memory}, nhờ đó
các ca kiểm thử chạy độc lập, không phụ thuộc vào SQL Server và không để lại dữ
liệu thừa. Mỗi ca kiểm thử tạo một cơ sở dữ liệu riêng mang tên ngẫu nhiên, bảo
đảm các ca không ảnh hưởng lẫn nhau.

Việc kiểm thử được ở mức này là kết quả trực tiếp của quyết định kiến trúc đã nêu
ở mục~\ref{sec:kien-truc}: vì logic nghiệp vụ nằm ở tầng Service chứ không nằm
trong Controller, nên có thể khởi tạo và kiểm thử trực tiếp mà không cần hạ tầng
web.

Cần lưu ý một giới hạn của cách tiếp cận này: nhà cung cấp In-Memory
\textbf{không thực thi} ràng buộc duy nhất, khóa ngoại hay quy tắc so sánh chuỗi
của SQL Server. Vì vậy các ca kiểm thử ở đây kiểm chứng logic của thuật toán, còn
các hành vi phụ thuộc vào cơ sở dữ liệu thật (tranh chấp ghi đồng thời, vi phạm
khóa ngoại, thứ tự sắp xếp theo bảng mã tiếng Việt) nằm ngoài vùng phủ và được
kiểm tra thủ công trên môi trường chạy thật.

\subsection{Kết quả chạy kiểm thử}

Bộ kiểm thử gồm \textbf{29 ca}, chia theo bốn nhóm ứng với ba thuật toán lõi và
chức năng tìm kiếm, sắp xếp. Danh sách đầy đủ từ TC-01 tới TC-29 kèm mục tiêu
kiểm chứng của từng ca được liệt kê ở Phụ lục F. Lệnh
\code{dotnet test src/CookingAdvisor.sln} cho kết quả:

\begin{codeblock}
Test run for .../CookingAdvisor.Tests.dll
    (.NETCoreApp,Version=v10.0)
A total of 1 test files matched the specified pattern.

Passed!  - Failed: 0, Passed: 29, Skipped: 0, Total: 29,
Duration: 337 ms - CookingAdvisor.Tests.dll (net10.0)
\end{codeblock}

\begin{table}[H]
\centering
\caption{Tổng hợp kết quả kiểm thử theo nhóm}
\label{tab:tong-hop-kiem-thu}
\begin{tabularx}{\textwidth}{|Y|P{2.0cm}|P{2.0cm}|P{2.4cm}|}
\hline
\bfseries Nhóm & \bfseries Số ca & \bfseries Đạt & \bfseries Không đạt \\ \hline
\code{SuggestionService} & 6 & 6 & 0 \\ \hline
\code{MenuPlannerService} & 9 & 9 & 0 \\ \hline
\code{ShoppingListService} & 8 & 8 & 0 \\ \hline
\code{RecipeService} & 6 & 6 & 0 \\ \hline
\bfseries Tổng & \bfseries 29 & \bfseries 29 & \bfseries 0 \\ \hline
\end{tabularx}
\end{table}

\subsection{Khiếm khuyết phát hiện qua kiểm thử}
\label{sec:khiem-khuyet}

Quá trình kiểm thử đã phát hiện hai khiếm khuyết mà quan sát bằng mắt thường khó
nhận ra. Đây là kết quả có giá trị nhất của phần kiểm thử.

\noindent\textbf{Khiếm khuyết 1: suy thoái tính đa dạng của thực đơn.}

\textit{Hiện tượng.} Khi thư viện món không đủ lấp 21 suất mà không lặp, thuật
toán chọn lại đúng một món cho hầu hết các suất còn trống.

\textit{Đo lường.} Trên bộ dữ liệu thử, một món xuất hiện 12 lần trong khi nhiều
món khác dùng được lại không được chọn lần nào. Con số này là kết quả đo tại thời
điểm phát hiện khiếm khuyết, trước khi vá.

\textit{Nguyên nhân.} Như đã phân tích ở mục~\ref{sec:thuat-toan-thuc-don}, trong
nhánh phải lặp, tiêu chí quyết định lùi về khoảng cách năng lượng. Do trạng thái
năng lượng tích lũy lặp lại theo chu kỳ, cùng một món luôn cho khoảng cách nhỏ
nhất.

\textit{Khắc phục.} Bổ sung tiêu chí ưu tiên món có số lần đã dùng ít nhất, đặt
trước tiêu chí năng lượng và chỉ áp dụng trong nhánh lặp.

\textit{Kết quả sau khắc phục.} Thực đơn sinh ra đạt \textbf{17 món khác nhau
trên 21 suất}. Hai ca kiểm thử hồi quy TC-10 và TC-11 được bổ sung để ngăn khiếm
khuyết tái diễn.

\noindent\textbf{Khiếm khuyết 2: thứ tự không xác định khi phân trang.}

\textit{Hiện tượng.} Khi sắp xếp theo năng lượng, các món có cùng giá trị năng
lượng không có thứ tự tất định giữa hai lần truy vấn, dẫn tới nguy cơ một món
xuất hiện ở cả trang trước và trang sau.

\textit{Nguyên nhân.} Câu lệnh sắp xếp chỉ có một khóa duy nhất. Với các bản ghi
có khóa bằng nhau, hệ quản trị cơ sở dữ liệu được tự do trả về theo thứ tự bất
kỳ.

\textit{Khắc phục.} Bổ sung khóa sắp xếp phụ là tên món cho mọi phương án sắp
xếp.

\textit{Kiểm chứng.} Ca kiểm thử TC-28 dựng 12 món có cùng giá trị năng lượng,
\textbf{nạp theo thứ tự tên giảm dần} để loại trừ khả năng kết quả đúng chỉ nhờ
trùng với thứ tự chèn dữ liệu, sau đó yêu cầu hai trang liên tiếp và kiểm tra
không có món nào lặp lại. Ca kiểm thử này đã được xác nhận là \textbf{có hiệu
lực}: khi tạm gỡ khóa sắp xếp phụ khỏi mã nguồn, ca kiểm thử chuyển sang trạng
thái không đạt; khi khôi phục, ca kiểm thử đạt trở lại.

\subsection{Kiểm thử giao diện}
\label{sec:kiem-thu-giao-dien}

Ngoài kiểm thử tự động ở mức đơn vị, giao diện được kiểm tra thủ công trên trình
duyệt theo các tiêu chí phi chức năng đã nêu ở
mục~\ref{sec:yeu-cau-phi-chuc-nang}.

\begin{table}[H]
\centering
\caption{Kết quả kiểm tra giao diện theo tiêu chí phi chức năng}
\label{tab:kiem-thu-giao-dien}
\begin{tabularx}{\textwidth}{|P{3.4cm}|Y|Y|}
\hline
\bfseries Tiêu chí & \bfseries Cách kiểm tra & \bfseries Kết quả \\ \hline
N4, không cuộn ngang & Đo \code{scrollWidth} so với \code{clientWidth} ở 375 px
  trên các trang chính & Đạt, không trang nào cuộn ngang \\ \hline
N5, tương phản WCAG AA & Đo tỉ lệ tương phản của các cặp màu chữ và nền
  trên trang thật & Đạt, cặp thấp nhất 4,69:1 so với ngưỡng 4,5:1 \\ \hline
Chế độ tối & Chuyển chế độ và đo lại toàn bộ cặp màu & Đạt ở cả hai chế độ \\ \hline
Điều hướng bàn phím & Kiểm tra bộ lọc dạng ngăn kéo: mở, giữ tiêu điểm,
  đóng bằng phím Esc & Đạt \\ \hline
Nhật ký lỗi trình duyệt & Theo dõi bảng điều khiển của trình duyệt khi duyệt các
  trang & Không có lỗi \\ \hline
\end{tabularx}
\end{table}

\section{Đánh giá}
\label{sec:danh-gia}

\subsection{Kết quả đạt được}
\label{sec:ket-qua-dat-duoc}

Đối chiếu với sáu mục tiêu đề ra ở phần Mở đầu:

\begin{table}[H]
\centering
\caption{Đối chiếu kết quả đạt được với mục tiêu đề ra}
\label{tab:doi-chieu-muc-tieu}
\begin{tabularx}{\textwidth}{|P{5.6cm}|Y|}
\hline
\bfseries Mục tiêu & \bfseries Kết quả \\ \hline
1. Quản lý kho dữ liệu món ăn & Đạt, 26 món với đầy đủ thuộc tính và định lượng
  nguyên liệu \\ \hline
2. Thuật toán gợi ý theo nguyên liệu & Đạt, 6 ca kiểm thử đều đạt \\ \hline
3. Thuật toán sinh thực đơn tuần & Đạt, 9 ca kiểm thử đều đạt \\ \hline
4. Sinh danh sách đi chợ & Đạt, 8 ca kiểm thử đều đạt \\ \hline
5. Xác thực và phân quyền & Đạt, kiểm chứng trên hệ thống đang chạy \\ \hline
6. Kiểm chứng bằng kiểm thử tự động & Đạt, 29 ca kiểm thử, phát hiện 2
  khiếm khuyết thực tế \\ \hline
\end{tabularx}
\end{table}

\subsection{Hạn chế}
\label{sec:han-che}

Hệ thống còn các hạn chế sau, được nêu đầy đủ để làm cơ sở cho phần Hướng phát
triển ở Chương~\ref{ch:ket-luan}:

\textbf{Một là gợi ý chưa xét định lượng:} thuật toán hiện chỉ xét việc có hay
không có một nguyên liệu, chưa xét người dùng có đủ khối lượng hay không, nên
một người có 50 gam thịt vẫn được gợi ý món cần 500 gam. \textbf{Hai là chưa cá
nhân hóa:} kết quả gợi ý chỉ phụ thuộc tập nguyên liệu đầu vào, giống nhau với
mọi người dùng, lịch sử nấu ăn chưa được khai thác. \textbf{Ba là dinh dưỡng ở
mức thô:} mới dừng ở năng lượng theo khẩu phần, chưa bóc tách đạm, béo, tinh bột
hay vi chất. \textbf{Bốn là} ảnh minh họa nhập bằng đường dẫn, chưa có chức năng
tải ảnh lên. \textbf{Năm là} chưa có đổi mật khẩu, quên mật khẩu và xác thực
email. \textbf{Sáu là} khu quản trị chưa phân trang và chưa có tìm kiếm, với quy
mô dữ liệu mẫu hiện tại thì chưa gây trở ngại nhưng sẽ là hạn chế khi số món
tăng lên. \textbf{Bảy là} thuật toán lập thực đơn theo hướng tham lam nên không
bảo đảm nghiệm tối ưu toàn cục về cân đối năng lượng. \textbf{Tám là} chưa có
kiểm thử tích hợp và kiểm thử giao diện tự động, phần giao diện mới được kiểm
tra thủ công.
